% -----------------------------
\section{THEORETICAL BACKGROUND}
%------------------------------

\textcite{campos2023} modeled a preloaded spring set used as part of the secondary suspension scheme of passenger coaches with Y32 bogies in Brazilian railroads, and composed of one coil spring in parallel with rubber springs. Data from a compression test was used to determine the rubber's hyperelastic material parameters according to the Mooney-Rivlin model. Finite elements analysis was used to obtain spring stiffness for each preload. The subsequent non-linear stiffness behavior was approximated by a piece-wise linear model. Nadal's ratio and wear number were calculated using a multibody analysis to compare the derailment risks predicted by a constant stiffness model, against those predicted by the piecewise model. Results suggest that employing exact material characterization of elastomeric components shows great potential towards good design and operation of coaches.

\textcite{kang2022} summarized the uses of rubber materials in railway vehicles, highlighting the importance of bushing elements to improve damping characteristics and enhance ride comfort. The authors present the constitutive theories of hyperelascity, and the relevant testing methods and standards for these materials. They present the mechanisms of aging and fatigue, and stress the importance of selecting correct fatigue criteria based on the actual fatigue failure index of components. Furthermore a need is identified, for a mathematical model capable of describing the failure characteristics of rubber components in different service conditions.

The mechanical behavior of a conical spring under multiaxial loads and pre-compression was obtained from single axis stiffness tests by \textcite{lalo2019}. The authors used the Hooke-Jeeves pattern seach optimization algorithm to fit the simulated data obtained through a hybrid finite elements formulation with uncoupled pressure terms. Six different hyperelastic models were used to characterize the spring. For each model, the error of the simulated FE analysis was evaluated against prototype test data. The optimization algorithm was then employed for the Arruda-Boyce model due to the observed superior correlation to experimental data. Comparison of simulated results against experimentally obtained load-deflection curves shows higher inaccuracy of measured material parameters for components displaying complex geometric configurations.

An inverse nonlinear finite elements method was developed by \textcite{lei2013} to iteratively reconstruct the stiffness information of a rubber bushing component with complex geometry from experimental data. For each iteration, the radial and axial stiffness curves were obtained via FEA, and compared to the experimental data. A genetic algorithm was used to optimize the weighted value of four optimization subobjectives, involving the six coefficients of the three-term Ogden hyperelastic model. 

\textcite{lima2022} investigated the dynamic behavior of shear pads, historically used on the primary suspension of freight vehicles, to better control warping and shearing of trucks. Building on the results of a compression and torsion test, they extracted a linear approximation of the pad's load-displacement curve. Friction contact was assumed on the surface of a 3D Abaqus\textregistered model of the pad, which was solved using the method of the augmented lagrangian. The Abaqus\textregistered internal fitting routine was used to obtain the Neo-Hookean hyperelastic characteristics of the pad from experimental data. A second FEM model was used to simulate the side frame used for pad support. The stiffness matrices from both models was then extracted, and used on a dynamic multibody model of a heavy-haul vehicle built on SIMPACK\textregistered. Results show that shear pads reduces wear number by 43\%. Multibody analysis also showed that wheelset lateral displacement was damped, while yaw angles were increased. Lastly, the longitudinal stiffness was found to be the only one relevant towards wear number reduction.

Another element which improves ride comfort is the suspension control arm. \textcite{tang2014} applied the SIMP methodology for the topology optimization of the control arm. Compatibility conditions were applied to the contact between control arm and joints, and the static response was obtained via FEM for two scenarios: with and without the effect of ball joint and rubber bushings stiffness. The authors conclude that consideration of stiffness of rubber bushings and ball joints is paramount to the static analysis, while the approach used can be used to model similar components where rubber bushings are present.

\textcite{vinith2024} investigated the performance of different rubber materials used for train suspension bushes undergone the process of vulcanization. A three-stage testing is employed to asses each component's response under compression, tensile strength, and fatigue conditions. The authors find that natural rubber exhibits superior fatigue resistance characteristics, along with better response under tensile and compression conditions relative to other candidates. Their findings highlight the relevance of thorough material selection, and serves as an important guide towards appropriate rubber material selection for suspension bushes.